\documentclass{article}

\usepackage{microtype}
\usepackage{graphicx}
\usepackage{booktabs}
\usepackage{hyperref}

\newcommand{\theHalgorithm}{\arabic{algorithm}}
\usepackage{icml2026}

\InputIfFileExists{values.tex}{}{}
% PROTECTED HARNESS METADATA. Agents use these macros; do not edit them.
\newcommand{\modelName}{Qwen2.5-0.5B-Instruct Q8}
\newcommand{\modelFile}{\texttt{qwen2.5-0.5b-instruct-q8\_0.gguf}}
\newcommand{\modelSha}{\texttt{ca59ca7f13d0e15a8cfa77bd17e65d24f6844b554a7b6c12e07a5f89ff76844e}}


\icmltitlerunning{\smash{Answer Position and a Lexical Shortcut}}

\begin{document}

\twocolumn[
\icmltitle{Where the Answer Sits and What a Filter Keeps:\
           A Synthetic Retrieval Diagnosis for a Small Language Model}
\icmlsetsymbol{equal}{*}
\begin{icmlauthorlist}
\icmlauthor{Anonymous Authors}{anon}
\end{icmlauthorlist}
\icmlaffiliation{anon}{Anonymous Institution}
\icmlcorrespondingauthor{Anonymous Author}{anon.email@domain.com}
\icmlkeywords{long context, retrieval, position effects, small language models}
\vskip 0.3in
]
\printAffiliationsAndNotice{}

\begin{abstract}
We test whether answer position and lexical competition change synthetic clause retrieval by a frozen small language model. The experiment uses \sampleSize{} deterministic probes with mean prompt lengths of \tierTwoTokens{}, \tierFourTokens{}, and \tierEightTokens{} tokens. Moving the same answer clause from the start to the end raised accuracy by \positionTwoStartEndChange{} points at the shortest tier and \positionFourStartEndChange{} points at the middle tier. A word-overlap filter retained a question-like decoy and reduced accuracy from \pruningFullAccuracy{}\% to \pruningPrunedAccuracy{}\%; removing only that decoy reversed the failure. An exact-name selector reached \entitySelectorAccuracy{}\% on its \entitySelectorCount{} template-matched cases, but deterministic extraction tests show that this result depends on a benchmark shortcut and does not establish a general mitigation. The matched length comparison remains unresolved, so the practical conclusion is limited to testing lexical filters against plausible competing clauses before deployment.
\end{abstract}

\section{Introduction}

You paste a long contract into a chatbot and ask, ``What code renews the vendor agreement?'' The answer sits in one clause, but a similar buyer clause repeats the words in your question and gives a different code.

This lookup can fail for distinct reasons. More text may burden the model, the answer may occupy a weak position, or a filter may preserve the wrong clause because its wording resembles the question. An experiment must hold the scenario fixed to distinguish these explanations.

We study a synthetic, gradeable version of the contract situation. Each document contains an answer clause for one company and a competing clause for another. The question shares its wording with the competitor, while the answer clause expresses the same relation in different words. This design follows NoLiMa's warning that literal question-to-answer overlap can make retrieval tests too easy \citep{modarressi2025nolima}, but it also creates a different hazard: entity repetition can reveal the answer sentence without model retrieval.

Our contribution is a controlled diagnosis on the frozen \modelName{} artifact. We move the same answer clause while fixing the competitor's final sentence index, compare full documents with word-overlap filtering, and remove only the retained competitor to locate the filter's failure. We also test the exact-name rule against aliases, pronouns, repeated or missing entities, ambiguity, and unseen question forms. These tests frame its perfect template result as a benchmark shortcut rather than a general context selector.

The experiment asks three questions:
\begin{enumerate}
    \item When all other document content stays fixed, how does moving the same answer clause among the start, middle, and end change retrieval accuracy?
    \item On matched cases, which sentences does word-overlap filtering retain, and does the retained competitor account for the model's errors after filtering?
    \item Across scenario families present at every tier, how much does accuracy change as mean prompt length grows from \tierTwoTokens{} to \tierEightTokens{} tokens?
\end{enumerate}

The observed families show strong end-position gains and a decoy-driven pruning failure. The evidence does not identify a prompt-length threshold or an isolated effect of filler relevance.

\section{Related Work}

\citet{liu2023lost} found stronger use of information near context boundaries than in the middle for multi-document question answering and key-value retrieval. We conduct a smaller clause-reordering test on a laptop-scale model. Our pattern chiefly favors recent information rather than reproducing the full U-shape.

Matched padding experiments show that performance can decline before a model reaches its context limit \citep{levy2024same}, and RULER finds that usable context can be shorter than the advertised window \citep{hsieh2024ruler}. Our sparse shared-family comparison cannot establish such a decline. Work on irrelevant statements also reports disrupted reasoning \citep{shi2023distracted}; our two filler generators change several properties together, so their contrast cannot isolate relevance.

NoLiMa reduces literal overlap between questions and answer-bearing text to expose retrieval difficulty \citep{modarressi2025nolima}. We add a paired failure diagnosis: surface overlap makes a simple filter preserve the competitor, while exact entity repetition makes another rule identify the answer sentence. The tested artifact belongs to the instruction-tuned Qwen family \citep{qwen2024qwen25}.

\section{Probe Design}

\paragraph{Model and execution.}
We evaluate \modelFile{} with the following SHA checksum:
\par\noindent\resizebox{\linewidth}{!}{\modelSha}\par
We use \texttt{llama-completion} with greedy decoding, zero temperature, and a fixed seed. The manifest contains \manifestProbeCount{} probes, each stored as one prompt and one raw output under a stable identifier. Interrupted inference resumes by skipping existing outputs.

\paragraph{Matched scenario families.}
Seeded code generates a target company and code, a competing company and code, a question, and filler text without using answer position in the seed. The question asks for the code required to extend the target company's contract. Its wording matches the competing clause, whereas the answer clause calls the value a renewal identifier. A literal control reverses these phrasings.

For each full-document family, we assign final sentence slots before placing either special clause. The answer then moves among start, middle, and end while the question, codes, clause text, filler order, and competitor's final sentence index remain fixed. Deterministic tests verify this invariant across every generated family and tier. The resulting contrasts isolate answer relocation within this synthetic document, although end placement remains diagnostic because a user would need to know the answer location before applying it.

\paragraph{Tokenizer-measured prompt amount and confounded filler bundles.}
The model tokenizer measures each prompt. Full-document tiers average \tierTwoTokens{}, \tierFourTokens{}, and \tierEightTokens{} tokens. Cases used for prompt-amount comparisons occur at every tier with the same family, answer position, question, and filler condition.

One filler generator produces business clauses containing other companies and code-shaped candidates; the other produces unrelated descriptions without those candidates. Thus, their contrast changes domain, entity and code density, candidate count, and sentence templates along with topical relation. An experiment identifying relevance would hold those properties fixed in a matched factorial design and change only whether filler bears on the question. We do not report the present bundle contrast as a relevance effect.

\paragraph{Filters and diagnostic controls.}
At the middle tier, the word-overlap filter keeps sentences sharing at least two non-stopword terms with the question. When fewer than \pruneFloor{} qualify, it keeps the \pruneFloor{} highest-scoring sentences, breaks ties by original order, and preserves document order. We compare only cases having both full and filtered outputs.

The decoy-removal control deletes the known competing clause from the filtered prompt while leaving all other retained sentences fixed. It uses a synthetic answer label and is diagnostic, not deployable. A separate exact-name rule parses the generator's fixed question form and keeps a sentence only when it is the unique exact mention of that company. On the generated templates, that sentence contains one code, so a deterministic rule can extract the code without invoking the language model. Unit tests reject aliases, pronouns, repeated names, absent names, multiple codes, and unseen question forms; failures use the word-overlap fallback in the inference pipeline. We therefore use the exact-name condition only to diagnose benchmark structure.

\section{Scoring and Comparisons}

A response is accurate when it contains the target code without the competing code. It complies exactly when it contains one code and no other text. A distractor capture contains the competitor without the target; an invalid response contains both expected codes or neither.

All estimates describe the observed fixed-seed families. Each display gives probe and family counts, and paired changes use only cases present in both conditions. Wins are cases corrected by the second condition, losses become incorrect, and ties do not change. The full-document tiers contain only \tierEightFamilyCount{} to \tierTwoFamilyCount{} families, so we make no population confidence claim. Greedy decoding removes sampling variation across identical reruns.

\section{Results}

\paragraph{Position sensitivity.}
The first question has a consistent descriptive answer for start versus end placement (Table~\ref{tab:position}). At the shortest tier, accuracy rose from \positionTwoStartAccuracy{}\% at the start to \positionTwoEndAccuracy{}\% at the end, a paired change of \positionTwoStartEndChange{} points across \positionTwoStartEndDenominator{} cases from \positionTwoStartEndFamilyCount{} families. The middle-tier change was \positionFourStartEndChange{} points across \positionFourStartEndDenominator{} cases from \positionFourStartEndFamilyCount{} families. The longest tier favored the end by \positionEightStartEndChange{} points, but only \positionEightStartEndFamilyCount{} families contributed.

\begin{table}[t]
\caption{Accuracy by answer position. Counts are matched probes/families; W/L/T compares start with end.}
\label{tab:position}
\centering
\small
\begin{tabular}{@{}lccccc@{}}
\toprule
Tokens & Start & Middle & End & Pairs/fam. & W/L/T \\
\midrule
\tierTwoTokens{} & \positionTwoStartAccuracy{} & \positionTwoMiddleAccuracy{} & \positionTwoEndAccuracy{}
& \positionTwoStartEndDenominator{}/\positionTwoStartEndFamilyCount{}
& \positionTwoStartEndWins{}/\positionTwoStartEndLosses{}/\positionTwoStartEndTies{} \\
\tierFourTokens{} & \positionFourStartAccuracy{} & \positionFourMiddleAccuracy{} & \positionFourEndAccuracy{}
& \positionFourStartEndDenominator{}/\positionFourStartEndFamilyCount{}
& \positionFourStartEndWins{}/\positionFourStartEndLosses{}/\positionFourStartEndTies{} \\
\tierEightTokens{} & \positionEightStartAccuracy{} & \positionEightMiddleAccuracy{} & \positionEightEndAccuracy{}
& \positionEightStartEndDenominator{}/\positionEightStartEndFamilyCount{}
& \positionEightStartEndWins{}/\positionEightStartEndLosses{}/\positionEightStartEndTies{} \\
\bottomrule
\end{tabular}
\end{table}

\paragraph{Lexical competition.}
The second question is answered by what the filter retained. Word-overlap filtering kept both the answer and competitor in every matched prompt, and accuracy fell from \pruningFullAccuracy{}\% to \pruningPrunedAccuracy{}\% (Table~\ref{tab:diagnosis}). Removing only the competitor raised accuracy to \pruningNoDecoyAccuracy{}\%, a paired change of \pruningNoDecoyChange{} points with \pruningNoDecoyWins{} wins and \pruningNoDecoyLosses{} losses. This intervention identifies the retained competitor as the cause of the filter's observed failure in this grid.

The exact-name condition reached \entitySelectorAccuracy{}\% on \entitySelectorCount{} cases from \entitySelectorFamilyCount{} families after retaining every answer clause and no competitor. That result measures the template invariant exposed by deterministic extraction, not robust language-model retrieval. Across all outputs, distractor capture was \distractorCaptureRate{}\%, exact-output compliance was \exactOutputComplianceRate{}\%, and invalid output was \invalidOutputRate{}\%.

\begin{table}[t]
\caption{Matched filter diagnosis. Counts are probes/families. The exact-name row is a benchmark-shortcut diagnostic.}
\label{tab:diagnosis}
\centering
\small
\begin{tabular}{@{}lcc@{}}
\toprule
Condition & Accuracy & Probes/families \\
\midrule
Full document & \pruningFullAccuracy{} & \pruningFullCount{}/\pruningFullFamilyCount{} \\
Word overlap & \pruningPrunedAccuracy{} & \pruningPrunedCount{}/\pruningPrunedFamilyCount{} \\
Decoy removed & \pruningNoDecoyAccuracy{} & \pruningNoDecoyCount{}/\pruningNoDecoyFamilyCount{} \\
Exact name & \entitySelectorAccuracy{} & \entitySelectorCount{}/\entitySelectorFamilyCount{} \\
\bottomrule
\end{tabular}
\end{table}

\paragraph{Prompt amount.}
The third question remains unresolved. Accuracy on shared cases was \amountSharedTwoAccuracy{}\% at \tierTwoTokens{} tokens and \amountSharedFourAccuracy{}\% at \tierFourTokens{} tokens. The paired change was \amountSharedTwoFourChange{} points, comprising \amountSharedTwoFourWins{} wins, \amountSharedTwoFourLosses{} losses, and \amountSharedTwoFourTies{} ties over \amountSharedTwoFourDenominator{} cases from \amountSharedTwoFourFamilyCount{} families. At \tierEightTokens{} tokens, accuracy was \amountSharedEightAccuracy{}\%; the intersection has only \amountSharedEightCount{} probes from \amountSharedEightFamilyCount{} families. These observations neither locate a degradation threshold nor show that length matters less than position.

\section{Limitations}

The evidence concerns one small quantized model, few fixed-seed scenario families, and synthetic clause lookup rather than natural contract analysis. Retrieval accuracy is not general capability, and the observed position pattern may not transfer to other models or documents. End placement is not a deployable remedy when the answer location is unknown. The sparse longest tier prevents a conclusion about prompt-length degradation.

The filler bundles cannot identify a relevance effect because they jointly change domain, entity and code density, candidate count, and templates. The exact-name condition is guaranteed unusually favorable structure: the fixed question names one company, one sentence uniquely repeats it, and that sentence contains one code. Its failure tests document important unsupported cases but are unit tests, not model evaluations on a broader probe distribution. The decoy-removal control also uses a known synthetic label unavailable in natural documents.

\section{Conclusion}

For this model and grid, moving a matched answer clause to the end improved retrieval, while word-overlap pruning strengthened a lexical competitor. Removing only that competitor reversed the pruning failure, which ties the error to selection rather than context reduction alone. Exact-name selection exposes a shortcut in the synthetic template and should not be read as a general mitigation. Builders using cheap filters for offline document question answering should test whether plausible decoys outrank the answer before treating shorter context as safer.

\section{Reproducibility}

Every reported quantity is generated from the committed manifest and raw model outputs. Following the principle that computational reports should expose budgets and variation \citep{dodge2019show}, we record the frozen model, decoding settings, stable probe identifiers, and the \cpuBudgetMinutes{}-minute CPU budget. From cached outputs, \texttt{make all} regenerates results, value macros, and the paper; \texttt{make clean-deep \&\& make all} repeats deterministic inference. Repository gates check frozen inputs, citations, numeric provenance, pipeline freshness, result sanity, prose style, anonymity, and compilation.

\label{main-body-end}
\bibliography{refs}
\bibliographystyle{icml2026}

\end{document}
